\documentclass[11pt,a4paper]{article}
\usepackage[T1]{fontenc}
\usepackage[utf8]{inputenc}
\usepackage{lmodern,microtype}
\usepackage[margin=25mm]{geometry}
\usepackage{amsmath,amssymb,amsthm}
\usepackage{booktabs,enumitem}
\usepackage[round,authoryear]{natbib}
\usepackage{xcolor}
\usepackage[colorlinks=true,allcolors=blue!45!black]{hyperref}
\hypersetup{pdftitle={Sharp overstatement of coalition savings in newsvendor relaxations},
  pdfauthor={Guido Morgante}}
\setlist{nosep,leftmargin=*}
\setlength{\emergencystretch}{2em}
\newtheorem{theorem}{Theorem}
\newtheorem{proposition}[theorem]{Proposition}
\newtheorem{corollary}[theorem]{Corollary}
\newtheorem{lemma}[theorem]{Lemma}
\theoremstyle{remark}
\newtheorem{remark}[theorem]{Remark}
\newcommand{\E}{\mathbb E}
\newcommand{\R}{\mathbb R}
\newcommand{\one}{\mathbf 1}
\newcommand{\pos}[1]{\left[#1\right]_+}
\newcommand{\subsets}{\mathcal S}
\newcommand{\gap}{\Delta}
\DeclareMathOperator{\Var}{Var}
\title{Sharp overstatement of coalition savings\\in newsvendor relaxations}
\author{Guido Morgante}
\date{Working draft --- 5 November 2026}

\begin{document}
\maketitle
\begin{abstract}
Fractional coalition optimization gives an upper bound on the largest saving
available to a coalition under a fixed cost allocation. We characterize its
worst-case overstatement in an exclusion hierarchy for a newsvendor cost
game. For every prescribed positive share vector, the supremum of the gap
divided by expected grand-offer cost is the sum of the largest retained
shares at each nonterminal exclusion depth. A common sequence of bounded
two-state instances approaches these values, has a correct forecast of total
output, and leaves every actual proper coalition strictly stable. The
construction also gives a sharp root bound and a rank-one norm-cost
representation. A standalone-cost condition reduces the bound; an iid
validation argument and a descriptive comparison delimit its practical scope.
The results concern the error of fixed-allocation fractional certification.
They do not establish general statistical or computational superiority.
\end{abstract}

\section{Introduction and relation to prior work}

Suppose a group implements a common newsvendor decision and divides its
expected cost according to fixed shares. A coalition can compare its assigned
charge with the cost of acting independently. Assessing the largest saving
requires a search over subsets, and allowing fractional memberships provides
a continuous upper bound. A positive fractional value can nevertheless be a
false indication of a profitable departure: the optimizing membership may
not describe any actual coalition.

This note asks a quantitative question. For a fixed positive cost split,
how large can that overstatement be after requiring several participants to
be completely absent? The answer is an exact worst-case constant, indexed
by the exclusion depth and the shares. The upper estimate follows from the
charge budget. The matching construction realizes that estimate
asymptotically while retaining positive producer variances, bounded outputs,
an aggregate-correct forecast and strict stability of every proper coalition.
The same sequence works at all nonterminal depths.

Fractional participation and comparisons between ordinary cores and dual
allocations are established ideas. \citet{aubin1981fuzzy} develops cooperative
fuzzy games; \citet{samet1984dualcore} studies LP-game cores, dual sets and
replication. \citet{denault2001coherent} connects homogeneous fractional cost
inequalities to subdifferentials. These results motivate the distinction
between integer and fractional stability, rather than making that
distinction a contribution of this note.

The closest optimization formulation is the fixed-allocation subproblem in
the 2005 manuscript of \citet[Sections 5.2--5.3, Eq.~19]{burer2005centralized}.
Section~\ref{sec:root} gives its exact correspondence with our root problem
for norm-valued costs. The manuscript subsequently optimizes the allocation;
our result keeps it fixed. That relaxation should be cited to the 2005
manuscript: the journal version \citep{burer2012newsvendor} instead develops
fair allocations under optimized Gaussian correlations and does not retain
the same Section~5.

Partially continuous coalition searches also appear in
\citet[Sections 4--6]{bonnans2008leastcore}. Their selected continuous
classes differ from the union of zero-coordinate faces used here. The
absolute-sum two-state cost in our example is a familiar device:
\citet[Example 1 and Theorem 7]{chenzhang2009inventorygames} distinguish core
and dual allocations and establish hardness of core membership testing.
We do not claim a new hardness result, general relaxation scheme or
two-state representation. The proposed increment is the exact extremal
value for every prescribed positive split and exclusion depth under the
stated realizability restrictions.

\section{Model and exclusion hierarchy}\label{sec:model}

Let $N=\{1,\ldots,n\}$, $n\ge3$, and let $Y=(Y_1,\ldots,Y_n)$ be jointly
distributed with $0\le Y_j\le R_j<\infty$. Within-vector dependence is
unrestricted. For positive penalties $\alpha,\beta$, use
\begin{equation}\label{eq:loss}
 \ell(q,t)=\alpha(t-q)_++\beta(q-t)_+.
\end{equation}
For memberships $z\in\R_+^n$, define the \emph{physical fractional cost}
\begin{equation}\label{eq:cost}
 c(z)=\min_{q\in\R}\E\ell(q,z^\top Y).
\end{equation}
An optimizer lies in $[0,R^\top z]$. Fix an implemented offer
$q_0\in[0,\sum_jR_j]$ and positive shares $w_j$ summing to one. Set
\[
 L=\E\ell(q_0,\one^\top Y),\qquad a_j=Lw_j,\qquad
 f(z)=a^\top z-c(z).
\]
Unless stated otherwise $L>0$. The choice $q_0$ need not be optimal under
the true law. Charges are frozen before the coalition search; they are
not reallocated after observing a candidate departure.

Write $a(S)=\sum_{j\in S}a_j$ and
$\subsets=\{S:\emptyset\ne S\subsetneq N\}$. The largest positive saving is
\begin{equation}\label{eq:B}
 B=\max\{0,\max_{S\in\subsets}f(\one_S)\}.
\end{equation}
We call a proper coalition strictly stable if $f(\one_S)<0$.
If $L=c(\one)$, the additional efficiency condition holds and $B=0$
means that $a$ is an ordinary-core allocation. Without this equality,
we use the excess interpretation and do not assert ordinary-core membership.

\begin{lemma}\label{lem:sublinear}
The function $c$ is nonnegative, convex, positively homogeneous and
subadditive on $\R_+^n$, with $c(0)=0$.
\end{lemma}
\begin{proof}
Nonnegativity is immediate. Scaling both $q$ and $z$ in
\eqref{eq:loss} proves positive homogeneity. Given two optimal offers,
their convex combination is an admissible offer at the corresponding
convex combination of memberships. Joint convexity of the loss therefore
proves convexity of the optimized value. Convexity and homogeneity give
$c(z+z')=2c((z+z')/2)\le c(z)+c(z')$.
\end{proof}

At the root, define
\[
 P_0=\{z\in[0,1]^n:\one^\top z\le n-1\},\qquad
 D_0=\max_{z\in P_0}f(z).
\]
For a set $J$ of excluded participants, let
$P_J=\{z\in[0,1]^n:z_j=0\ (j\in J)\}$. For $1\le k\le n-1$, put
\begin{equation}\label{eq:Dk}
 D_k=\max\left\{0,\
    \max_{\substack{S\in\subsets\\|N\setminus S|<k}}f(\one_S),\
    \max_{\substack{J\subset N\\|J|=k}}\max_{z\in P_J}f(z)\right\}.
\end{equation}
An empty maximum is ignored. The direct integer checks retain coalitions
larger than a depth-$k$ face; omitting them would change the problem.

\begin{proposition}\label{prop:hierarchy}
$D_0\ge D_1\ge\cdots\ge D_{n-1}=B$.
\end{proposition}
\begin{proof}
Every proper integer coalition is either checked directly or belongs to
one of the faces at each depth. Every depth-$k+1$ face belongs to a
depth-$k$ face, as do the extra integer coalitions introduced at depth
$k+1$. All depth-one faces belong to $P_0$. At depth $n-1$, each face
has one free coordinate, and homogeneity makes its objective linear along
that segment. An endpoint attains the maximum. Together with the direct
checks these endpoints recover all proper integer coalitions and zero.
\end{proof}

For an explicit finite distribution, the face search is a linear program:
introduce an offer and positive/negative imbalance variables for each
scenario in \eqref{eq:cost}. At depth $k$ it requires $\binom nk$
face problems and $\sum_{j=1}^{k-1}\binom nj$ direct checks. This is a
description of the hierarchy, not a polynomial-time claim when $k$ grows
with $n$.

\section{Sharp error for prescribed shares}\label{sec:sharp}

Let $s_k(w)$ be the sum of the $n-k$ largest entries of $w$.

\begin{theorem}\label{thm:sharp}
Fix $n\ge3$ and $w>0$ with $\sum_jw_j=1$. For $1\le k\le n-2$,
\begin{equation}\label{eq:sharp}
 0\le D_k-B\le\pos{Ls_k(w)-B},\qquad
 \sup\frac{D_k-B}{L}=s_k(w).
\end{equation}
The supremum is over bounded instances with these shares and $L>0$.
It is unchanged if one requires unit capacities, positive individual
variances, two equiprobable states, unit absolute loss, an optimal true
grand offer, a forecast with the correct aggregate law that induces the
prescribed shares, and strict stability of every proper coalition.
At $k=n-1$ the error is zero.
\end{theorem}

\begin{proof}
\textit{Upper bound.}
On a face with $n-k$ free coordinates, nonnegative cost gives
$f(z)\le a^\top z\le Ls_k(w)$. Every direct integer candidate is at most
$B$, so $D_k\le\max\{B,Ls_k(w)\}$. Proposition~\ref{prop:hierarchy}
supplies the lower bound and terminal exactness.

\textit{Bounded construction and forecast realization.}
Relabel so that $w_1$ and $w_n$ are the two largest shares, with
$w_1\ge w_n$; order the middle shares increasingly. For $m\ge2$ set
$H=m+2(n-2)$. First work in unscaled units, with two equiprobable states
\begin{equation}\label{eq:states}
 y^+=(m,2,\ldots,2,0),\qquad
 y^-=(0,\ldots,0,H-1).
\end{equation}
Their totals are $H$ and $H-1$. Under unit absolute loss take
$q_0=H-\tfrac12$, an optimal true offer, giving $L=\tfrac12$.
Define a forecast by $\widetilde Y_j=w_jT$, where $T$ has these same two
total values. Its aggregate law is correct. With score $\nu=+1$ at $T=H$
and $\nu=-1$ at $T=H-1$, one has $\E\nu=0$ and
\[
 \E[\widetilde Y_j\nu]=\frac{w_j}{2},\qquad
 \sum_j\E[\widetilde Y_j\nu]=\frac12.
\]
Thus the forecast's optimal dual allocation, normalized by its total
cost, gives exactly $w$. Dividing all states and offers by $H$ puts
both true and forecast outputs in $[0,1]^n$, keeps every true variance
positive, and changes $L$ to $1/(2H)$. All ratios of excess to $L$
are unchanged.

\textit{Strict stability of actual coalitions.}
In unscaled units let
$d=y^+-y^-=(m,2,\ldots,2,-(H-1))$. The cost is
\begin{equation}\label{eq:absolute}
 c(z)=\frac12|d^\top z|.
\end{equation}
Indeed, the minimum mean absolute deviation of two equiprobable
values is half their distance. A nonempty coalition not containing
$n$ has state difference at least two. A proper coalition containing
$n$ omits positive differences of total at least two; since
$d^\top\one=1$, its difference is at most $-1$.
Consequently $c(\one_S)\ge\tfrac12=L>a(S)$ for all
$S\in\subsets$. Therefore $B=0$, with strict inequalities.

\textit{Fractional witness and sharpness.}
Omit the $k$ smallest middle shares, include every other
positive-difference participant fully, and give participant $n$
membership
\[
 r=\frac{H-2k}{H-1}=1-\frac{2k-1}{H-1}.
\]
For $1\le k\le n-2$ this lies in $(0,1)$. The two weighted
production totals coincide, so the membership has zero cost.
Its objective divided by $L$ is
\begin{equation}\label{eq:family}
 s_k(w)-w_n\frac{2k-1}{H-1}.
\end{equation}
As $m\to\infty$, this tends to $s_k(w)$. The same sequence
works for every nonterminal $k$, proving the claimed supremum.
\end{proof}

For equal shares the nonterminal supremum is $(n-k)/n$.
In particular, for large $n$ a fixed number of exclusions need not give
a small error relative to expected cost, even when all actual coalitions
are strictly stable.

\begin{proposition}\label{prop:finite}
For every finite member of \eqref{eq:states}, expression
\eqref{eq:family} equals $D_k/L$, not merely a feasible lower bound.
\end{proposition}
\begin{proof}
Direct integer values are negative. On a face omitting $n$,
$f(z)/L=\sum_{i<n}(w_i-d_i)z_i\le0$, since $d_i\ge2>w_i$.
On a face containing $n$, fix the other memberships and write
$u=\sum_{i<n}d_iz_i$. At least one positive-difference participant
is absent, hence $0\le u\le H-2$. As a function of $z_n$, the
normalized objective has slopes $w_n+(H-1)>0$ and
$w_n-(H-1)<0$ on either side of $u/(H-1)$.
It is therefore maximized by canceling the difference:
$z_n=u/(H-1)\in[0,1]$.

After cancellation, each free positive coordinate has coefficient
$w_i+w_nd_i/(H-1)>0$ and is set to one. The best omissions minimize
the sum of these coefficients. Participant 1 has a coefficient at
least as large as every middle participant because $m\ge2$ and
$w_1\ge w_i$. Among the middle participants the added coefficient
is the same. Thus omitting the $k$ smallest middle shares is optimal,
giving \eqref{eq:family}.
\end{proof}

\begin{remark}[Normalization]\label{rem:scale}
The sharpness denominator is $L$, not the worst possible loss of the
implemented offer over the declared box,
\[
 A=\max\{\beta q_0,\alpha(\textstyle\sum_jR_j-q_0)\}.
\]
In the unit-capacity version of \eqref{eq:states},
$q_0=1-1/(2H)$ and $A=n-1+1/(2H)$. Hence
\[
 \frac{L}{A}=\frac{1}{2H(n-1)+1}\longrightarrow0.
\]
The limiting construction does not establish a large error relative to
$A$ or explain a finite-data certificate's failure at a fixed
percentage-of-$A$ precision target.
\end{remark}

\section{Root relaxation and the choice of extension}\label{sec:root}

\subsection{Correspondence with the earlier norm-cost formulation}

Put $P_1=\{z\in[0,1]^n:1\le\one^\top z\le n-1\}$.
Positive homogeneity gives
\begin{equation}\label{eq:rootmap}
 D_0=\max\{0,\max_{z\in P_1}f(z)\}.
\end{equation}
To see this, a positive objective at a point of mass $s<1$ increases
when $z$ is replaced by $z/s$. Each coordinate remains at most one
because $z_j\le s$; the new mass is one. Nonpositive objectives do
not improve on zero.

For $c(z)=\|Mz\|_2$, Eq.~19 of \citet{burer2005centralized} has value
$v(a)=\min_{z\in P_1}\{c(z)-a^\top z\}$. Thus
$D_0=[-v(a)]_+$. The analogous binary program maps to $B$.
This identity concerns a fixed $a$, and does not assess the quality
of an allocation selected by an outer optimization.

\begin{corollary}\label{cor:root}
For the model and prescribed shares of Theorem~\ref{thm:sharp},
\[
 0\le D_0-B\le\pos{L(1-\min_jw_j)-B},\qquad
 \sup\frac{D_0-B}{L}=1-\min_jw_j.
\]
\end{corollary}
\begin{proof}
$D_0\le\max_{z\in P_0}a^\top z=L(1-\min_jw_j)$.
Since $D_0\ge D_1$, the depth-one construction in
Theorem~\ref{thm:sharp} proves the matching supremum.
\end{proof}

The same cost obstruction belongs to the rank-one norm class:
\eqref{eq:absolute} is $\|Mz\|_2$ with $M=d^\top/2$.
For a Gaussian realization, let $Z$ be standard normal,
$\eta=\E|Z|=\sqrt{2/\pi}$, and $X=\mu+dZ/(2\eta)$.
The optimized absolute-loss cost of $z^\top X$ is
$|d^\top z|/2$. Taking $\mu\ge0$ also makes the optimal offer
$\mu^\top z$ nonnegative. The covariance is singular and Gaussian
outputs are not bounded or almost surely nonnegative. Those physical
properties are supplied by the separate two-state realization, not by
the Gaussian one.

\subsection{Canonical and physical fractional costs}

The canonical concave cover is an established construction
\citep{samet1984dualcore,shellshear2009fuzzy}. In cost signs it reads
\begin{equation}\label{eq:cover}
 \bar c(z)=\min_{\lambda_S\ge0}
 \left\{\sum_{\emptyset\ne S\subseteq N}\lambda_Sc(\one_S):
       \sum_S\lambda_S\one_S=z\right\}.
\end{equation}
This is not generally the physical cost \eqref{eq:cost}.

\begin{proposition}\label{prop:cover}
For the sublinear cost $c$,
$c(z)\le\bar c(z)$ and $\bar c(\one_S)=c(\one_S)$ for every
actual coalition. If $a$ is an ordinary-core allocation, then
$a^\top z\le\bar c(z)$ for every nonnegative membership.
\end{proposition}
\begin{proof}
Subadditivity and homogeneity bound $c(z)$ by every feasible
objective in \eqref{eq:cover}. At $z=\one_S$, using $S$ alone
proves equality. For a core allocation, multiply
$a(S)\le c(\one_S)$ by $\lambda_S$ and sum, then minimize.
\end{proof}

For example, use unscaled states $(2,2,0)$ and $(0,0,3)$ and equal
charges $a_j=1/6$. Every actual proper coalition is strictly stable.
At $z=(1,0,2/3)$, however,
\[
 c(z)=0,\qquad a^\top z=5/18,\qquad \bar c(z)=2/3.
\]
For the last equality, only $\{1\}$, $\{3\}$ and $\{1,3\}$ can be
used in \eqref{eq:cover}. If the pair coefficient is $t\le2/3$,
the cost is $2-2t$, minimized at $t=2/3$. Thus ordinary-core
preservation for a canonical extension does not imply exactness
of the physical fractional audit.

\section{A condition that reduces the error}\label{sec:rho}

Let $c_j=c(\one_{\{j\}})$ and define
\[
 \rho=\min\bigl(\{1\}\cup\{a_j/c_j:c_j>0\}\bigr).
\]
In particular, $\rho=1$ when all singleton costs vanish.

\begin{proposition}\label{prop:rho}
For $1\le k\le n-2$,
\begin{equation}\label{eq:rho}
 D_k-B\le(1-\rho)\pos{Ls_k(w)-B}.
\end{equation}
\end{proposition}
\begin{proof}
The case $\rho=0$ is Theorem~\ref{thm:sharp}.
On a depth-$k$ face with free set $F$ put
$x=\sum_{j\in F}a_j(1-z_j)$. Nonnegative cost gives
$f(z)-B\le a(F)-B-x$. Also,
\[
 c(\one_F)\le c(z)+\sum_{j\in F}(1-z_j)c_j
            \le c(z)+x/\rho.
\]
As $F$ is a proper actual coalition, $B\ge a(F)-c(\one_F)$.
It follows that $f(z)-B\le(1/\rho-1)x$.
The smaller of the two bounds is at most their weighted average
with weights $1-\rho$ and $\rho$, which cancels $x$ and leaves
$(1-\rho)(a(F)-B)$. Take the maximum over faces, zero and the
direct integer candidates, all of whose excess above $B$ is
nonpositive. Use $a(F)\le Ls_k(w)$.
\end{proof}

For example, covering at least half of every singleton's optimized
cost halves the share-only ceiling. This is an elementary sufficient
condition from subadditivity, not an independent priority claim.
The root is not included in \eqref{eq:rho}; its maximizer can have
all coordinates positive. If $L=0$, all positive excesses and
all $D_k$ are zero directly.

Another sufficient condition is symmetry. Under exchangeable
production and equal charges $a_j=a_0$, fix a face $F$ with
$|F|=n-1$. Averaging memberships over its permutations, then
using convexity and homogeneity, gives
$c(z)\ge(\sum_{j\in F}z_j)c(\one_F)/(n-1)$.
Its maximum excess is therefore attained at zero or $\one_F$,
so $D_1=B=[(n-1)a_0-c(\one_F)]_+$. Equal marginal laws alone
do not imply this exchangeability assumption.

\section{Numerical checks and scope of validation}\label{sec:evidence}

The calculations support implementation checks, not proof or
bibliographic priority. The complete deterministic grid uses
$n\in\{3,6,8\}$, $m\in\{2,20,200\}$ and two share vectors:
uniform shares and shares proportional to
$(n,1,2,\ldots,n-2,n-1)$. All 18 laws have strictly stable actual
proper coalitions. Exact enumeration and the face implementation
agree with \eqref{eq:family} at all 84 depth values, including
the terminal levels. Table~\ref{tab:sharp} shows the uniform
six-participant slice.

\begin{table}[htbp]
\centering
\begin{tabular}{lrrrrr}
\toprule
$m$ & $k=1$ & $k=2$ & $k=3$ & $k=4$ & $k=5$ \\
\midrule
2 & 0.814815 & 0.611111 & 0.407407 & 0.203704 & 0.000000 \\
20 & 0.827160 & 0.648148 & 0.469136 & 0.290123 & 0.000000 \\
200 & 0.832528 & 0.664251 & 0.495974 & 0.327697 & 0.000000 \\
\midrule
Supremum & $5/6$ & $2/3$ & $1/2$ & $1/3$ & $0$ \\
\bottomrule
\end{tabular}

\caption{Checked values of $D_k/L$ with $B=0$, six uniform shares
and $H=m+8$. The last row is the proved supremum; terminal depth
is exact. Values are generated from the saved deterministic results.}
\label{tab:sharp}
\end{table}

Two separate primal LPs also agree with the existing root dual
implementation on the same 18 laws. The rank-one and physical costs
agree at 1,968 checked membership points, and a separate cover LP
verifies the three-participant example in
Section~\ref{sec:root}. These checks do not assess a reoptimized
allocation or claim coverage from a finite test count.

Appendix~\ref{app:validation} turns \eqref{eq:rho} into an iid
validation bound using established statistical tools. A prespecified
development comparison used 96 positive-share laws and 12
exchangeable controls. Four original cases with nonpositive
shares were excluded before sampling. One independent pilot and
validation stream per law supplied prefixes of
$M\in\{256,1024,4096\}$ observations each. The fixed rule was
independent of both streams. Each of the 18 method--size combinations
used $\delta=0.05/18$, yielding simultaneous validity across
the 18 comparisons within a law, not joint coverage across all
108 laws.

The alternatives are a Bonferroni-calibrated version of the
paired ratio procedure; a marginal DKW bound using both halves
for fixed-rule moments; a capacity bound using both halves;
a preceding bound from pilot-fitted dual prices on faces; and
paired SAA blocks of 64 observations. The latter specialize the
optimality-gap construction of \citet{mak1999montecarlo}, with
bounded finite-sample calibration. Each candidate-dependent
witness is evaluated only on validation data. Averages and
counts describe this development grid, not repeated-sampling
power or a controlled runtime comparison.

\begin{table}[htbp]
\centering
\begin{tabular}{lrr}
\toprule
Method & Mean $G/A$ & $G\le0.05A$ \\
\midrule
Paired ratio (IUT/fixed sequence) & 0.176440 & 5/47 \\
Bonferroni ratio & 0.177906 & 4/47 \\
Marginal DKW, pooled moments & 0.170268 & 1/47 \\
Capacity bound, pooled moments & 0.201686 & 0/47 \\
Pilot-fitted face prices & 0.143603 & 0/47 \\
Paired SAA blocks (size 64) & 0.217701 & 0/47 \\
\bottomrule
\end{tabular}

\caption{The 47 heterogeneous positive-share laws at
$M=4096$ pilot plus $M=4096$ validation observations.
The estimand is $D_1-B$. Here $G$ is its upper confidence bound
and $A$ is the maximum declared grand-offer loss.
All six prespecified methods are shown.}
\label{tab:validation}
\end{table}

Among the 35 original laws with true $\rho\ge0.6$, the paired,
Bonferroni and DKW procedures all certified $\rho\ge0.5$.
The broader target, a certified gap at most $0.05A$ in at least
half the heterogeneous laws, failed: the paired procedure passed
5 of 47. Even exact population values in \eqref{eq:rho} meet
that target in only 8 of 47. Thus improving scalar estimation
alone cannot eliminate this corollary's population limitation.

Post-run combinations of the paired and face bounds yield a
mean normalized upper gap of $0.128273$, versus $0.129887$
for DKW plus face, a descriptive difference of about $1.24\%$.
These combinations do not establish general superiority.
The observed run had no gap-coverage failures, but the probability
guarantees depend on the assumptions and proof. The SAA baseline
used 9,072 blocks and 50,232 LP solves across the study; its
fixed block choice does not settle the bias--concentration tradeoff.

\section{Discussion}

The exact error constant depends only on the retained shares,
but its sharpness is realized by a specialized dependence
structure. The result warns against interpreting a large
fractional objective as evidence of a profitable actual departure,
even with an aggregate-correct forecast. Whole-participant
exclusions reduce the ceiling but need not make it small relative
to expected cost.

Several limits remain essential. The scale issue in
Remark~\ref{rem:scale} prevents an interpretation as a
fixed maximum-loss error. Correct aggregate forecasting permits
individual forecast errors. The bound concerns this particular
physical relaxation hierarchy and does not exclude better
certificates or algorithms. The standalone condition is useful
only on a restricted part of the tested grid. No guarantee for
dependent time-series validation, real-data forecasting advantage
or general runtime superiority follows.

The focused source comparison did not locate the exact
depth/share formula among the inspected results. This is not
a certificate of global priority. The intended contribution
is the quantitative extremal calculation with its realizability
properties; established optimization and statistical ingredients
are credited explicitly. The present document is a working
draft for assessment of that contribution's significance.

\appendix
\section{A finite-data conversion of the standalone condition}
\label{app:validation}

Assume bounded iid validation observations independent of a pilot
and of the fixed rule $(q_0,w)$. Within-vector dependence remains
arbitrary. Condition on the pilot and fit offers $q_j\in[0,R_j]$.
For a predetermined threshold $t\in(0,1]$, use paired variables
\[
 X_{j,t}=w_j\ell(q_0,Y(N))-t\ell(q_j,Y_j).
\]
Write $b_j=\E\ell(q_j,Y_j)\ge c_j$.
Let $\ell_{j,t}$ be a scalar lower confidence bound satisfying
\begin{equation}\label{eq:scalar}
 \Pr\{\ell_{j,t}>\E X_{j,t}\mid\text{pilot}\}\le\eta.
\end{equation}
These pointwise bounds need not hold jointly.

\begin{lemma}\label{lem:iut}
For a fixed increasing grid $0<t_1<\cdots<t_K\le1$, test a
threshold only if all $\ell_{j,t}>0$. Stop at the first failed
threshold and let $\widehat\rho_{\rm grid}$ be the last passed
threshold, or zero. Then
$\Pr\{\widehat\rho_{\rm grid}>\rho\mid\text{pilot}\}\le\eta$.
\end{lemma}
\begin{proof}
If $t>\rho$, some $j$ has $c_j>0$ and
$a_j<t c_j\le t b_j$. Hence $\E X_{j,t}<0$, and passing
the threshold implies the event in \eqref{eq:scalar} for this
one index. For the sequential procedure, overstatement of
$\rho$ requires passing the first fixed grid point above
$\rho$, if any. Its failure probability is at most $\eta$.
\end{proof}

The reasoning is the intersection--union and fixed-sequence
principle \citep{berger1996intersectionunion,hsu1999fixedsequence}.
It removes the need for producer/grid multiplicity factors
for this single global claim. It does not give simultaneous
coverage of all scalar intervals or dimension-independent power.

For completeness, the implemented scalar bounds combine
Hoeffding and empirical Bernstein inequalities
\citep[Theorem 4]{maurer2009empiricalbernstein}. For a variable
in an interval of width $W$, $M\ge2$ observations and unbiased
sample variance $v$, use
\[
 r_\eta=\min\left\{
 W\sqrt{\frac{\log(2/\eta)}{2M}},\
 \sqrt{\frac{2v\log(4/\eta)}{M}}
 +\frac{7W\log(4/\eta)}{3(M-1)}\right\}.
\]
Each inequality receives tail $\eta/2$; a union bound validates
their minimum. The lower bound is the sample mean minus $r_\eta$,
clipped to the known range. Here
$X_{j,t}\in[-tA_j,w_jA]$, where
$A_j=\max\{\beta q_j,\alpha(R_j-q_j)\}$.
For a two-sided interval, split its budget again between tails.

A pilot-selected coalition $S\in\subsets$, with offer
$q_S\in[0,R(S)]$, has feasible excess
\[
 e_S=w(S)L-\E\ell(q_S,Y(S))\le B.
\]
A scalar lower confidence bound, clipped below at zero, gives
$W_-\le B$. The empty witness instead gives $W_-=0$ exactly.
Obtain a two-sided grand-cost interval $[L_-,L_+]$ and split
the total budget $\delta$ into three equal parts: the grid,
the grand interval and this witness.

The deterministic bound
$c_j\le C_j=\alpha\beta R_j/(\alpha+\beta)$ follows by
using the offer $\alpha R_j/(\alpha+\beta)$ and comparing
the two endpoint losses. Define
\[
 \widehat\rho_{\rm cap}
 =\min\bigl(\{1\}\cup\{w_jL_-/C_j:C_j>0\}\bigr),\qquad
 \widehat\rho=\max\{\widehat\rho_{\rm grid},
                    \widehat\rho_{\rm cap}\}.
\]
The grand interval is clipped to $[0,A]$.

\begin{proposition}\label{prop:validation}
With probability at least $1-\delta$, simultaneously for
$1\le k\le n-2$,
\[
 \widehat\rho\le\rho,\qquad
 D_k-B\le (1-\widehat\rho)\pos{s_k(w)L_+-W_-}.
\]
The terminal-depth error is zero.
\end{proposition}
\begin{proof}
Lemma~\ref{lem:iut}, the grand interval and witness bound
hold together by a three-event union bound. On this event
the cap bound also holds, as it uses only $L_-$.
Apply Proposition~\ref{prop:rho} and monotonicity in
$\rho,L,B$. The same event proves every depth bound.
\end{proof}

The implementation uses the grid $0.05,0.10,\ldots,1.00$.
The pilot selects from leave-one-out coalitions and their
fitted offers, or an empty witness. A numerical inconsistency
triggers the safe fallback $\widehat\rho=0$, $W_-=0$,
$L_+=A$. This fallback is recorded and did not occur in
the saved comparison. This appendix is a conversion of an
elementary population condition using established methods.

\section{Reproducibility and review status}\label{app:reproduce}

The accompanying repository contains the full proofs underlying
earlier diagnostics, source-by-source reading depth, protocols,
negative results and source hashes. The manuscript tables are
generated from existing JSON outcomes by
\texttt{paper/build\_tables.py}; no new sampling is used to
assemble them. The deterministic exclusion family and root
correspondence are regenerated by
\texttt{scripts/check\_exclusion\_ceiling.py} and
\texttt{scripts/check\_exclusion\_prior\_art.py}.

The source comparison distinguishes the 2005 manuscript from
the later journal version and records partial reading where
full originals were not obtained. Tests and this draft's
internal review are not independent mathematical review.
No submission or external-review outcome is claimed.

\bibliographystyle{plainnat}
\bibliography{references}
\end{document}
